% Shared by main.tex and extended/main.tex. Lines marked % CHECK-DATA depend on the data.
\section{Limitations and conclusion}
\label{sec:limitations}

Every parameter is an LLM's belief. The \voiNMembers{} members come from \voiNMembers{} developers, none of them the developer of the fact-extracting model, and all read the same sourced facts; still, the parameters are not validated against human experts, LLM intervals are often too narrow~\citep{hobor2026bayesian}, and no evaluation publishes a sensitivity and specificity against real harm. The extracted facts are a second weak point: an agent chose which sentences of the sources to show.\ifextended{} \ifdefined\voinoctxEtaStarPSQMed With no facts shown at all, $\PS$ by $\etast$ is \voinoctxEtaStarPSQMed{} (against \voiEtaStarPSQMed{} with facts) and the stakes ratio \voinoctxStakesRatio{} (against \voiStakesRatio) (Section~\ref{sec:noctx}).\fi\fi{} % CHECK-DATA
Each decision is written from the evaluation's own sources, so the comparison is between the decisions as much as between the evaluations; pricing both groups against one shared decision would separate the two. Published physical-AI safety evaluations are few and uneven, so the result describes today's evaluations, not physical AI as such.

\label{sec:conclusion}
Under our model, physical-AI safety evaluations give an undecided developer less value per dollar than LLM safety evaluations because far less is at stake per release. In both groups, the evaluations that can change a decision repay their build quickly. Pricing evaluations by their value of information gives a budget rule that takes measured sensitivities, specificities and costs as they become available. Code and data: \anonlink. % CHECK-DATA
